\section{Minimum cyclic cuts can avoid every short cycle}
\label{sec:random-long-minimum-cut}

A minimum cyclic cut need not contain an edge on a short cycle.
The obstruction persists at any fixed cyclic connectivity.
This concerns a proposed proof strategy, not the sharp cycle-count conjecture.

\begin{lemma}[Joining poles with distant ports]
\label{lem:random-distant-port-gluing}
Let $X,Y$ be connected normalized $k$-poles, where $k\ge2$.
Assume that each has rank at least $k$, and that distinct ports in each pole have distance at least $2k$.
Join their corresponding ports by $k$ edges, obtaining a cubic graph $G$.
\begin{enumerate}
\item If both poles have no internal cyclic cut of size below $k$, then $\lambda_c(G)=k$.
\item If both poles have no internal cyclic cut of size at most $k$, then the joining cut is the unique cyclic cut of size at most $k$.
\end{enumerate}
Uniqueness identifies a cut with its unordered pair of shores.
\end{lemma}

\begin{proof}
Consider a cyclic cut $\delta_G S$ of size $d$.
For the first assertion suppose $d<k$; for the second suppose $d\le k$.
Its restriction to either pole has internal boundary at most $d$.
By the corresponding connectivity hypothesis, the two restricted shores cannot both contain cycles.
They cannot both be forests either.
Indeed, if their component counts are $c_1,c_2$, then the pole rank would be
\[
 |\delta_X(S\cap X)|-c_1-c_2+1
 \le d-1<k.
\]
The empty-shore cases also exclude two forest sides, since the whole pole is cyclic.
Thus each restricted partition has a unique acyclic side.

Every component $B$ of such an acyclic side contains at most one port.
Otherwise a path in $B$ between two ports has length at least $2k$, so $|B|\ge2k+1$.
If $t_B$ counts its ports, the degree identity gives
\[
 |\delta_X B|=|B|+2-t_B\ge2k+3-k>k.
\]
This contradicts $|\delta_X B|\le d$.
In particular, an acyclic side containing $\alpha$ ports has internal boundary at least $2\alpha$.

Call the other restricted side the cyclic side.
If both poles' cyclic sides lie on the same side of $\delta_G S$, the opposite side is a forest.
Each of its tree components before the joining edges are restored has at most one port.
The joining edges can therefore only join two distinct tree components in pairs.
They cannot create a cycle.
This contradicts the cyclicity of the chosen cut.

The cyclic sides must therefore lie on opposite sides.
Let $\alpha,\beta$ count the ports in the two acyclic sides.
Relative to the joining cut, moving these ports can remove at most $\alpha+\beta$ crossing joining edges.
The new internal boundaries contribute at least $2\alpha+2\beta$.
Consequently
\[
 d\ge k-\alpha-\beta+2\alpha+2\beta
   =k+\alpha+\beta.
\]
This excludes $d<k$.
If $d\le k$, equality forces $\alpha=\beta=0$ and both internal boundaries to vanish.
Connectedness of the poles then makes both acyclic sides empty.
The cut is exactly the joining cut.
\end{proof}

\begin{theorem}[Long cycles through every minimum-cut edge]
\label{thm:random-long-minimum-cut}
For every fixed integer $k\ge5$, there are a constant $c_k>0$ and arbitrarily large connected simple cubic graphs $G$ such that
\[
 \lambda_c(G)=k,\qquad \operatorname{girth}(G)\ge k+1,
\]
the minimum cyclic cut is unique, and every cycle containing any edge of that cut has length at least
\[
 c_k\,\beta(G).
\]
\end{theorem}

\begin{proof}
Put $K=k+1$.
Choose a sufficiently large simple cubic graph $H$ with cyclic connectivity at least $K$ and girth at least $K$.
Such graphs exist at arbitrarily large orders.
For example, Wormald's random-regular connectivity theorem supplies them for every fixed $K$:
\href{https://doi.org/10.1016/S0095-8956(81)80021-4}{N.~C.~Wormald, \emph{The asymptotic connectivity of labelled regular graphs}, JCTB 31 (1981), 156--167}.
The theorem states that, asymptotically, random regular graphs conditioned on the required fixed girth have the corresponding cyclic connectivity.

Choose $2K+1$ edges of $H$ whose endpoint sets have pairwise distance at least $8K$.
A greedy choice is possible when $H$ is sufficiently large, since cubic balls of fixed radius have bounded size.
Divide these edges into a left set and a right set of size $K$, and one extra edge.
Also choose $H$ large enough that $\beta(H)\ge K$.
Write $n_0=|V(H)|$.

Construct a chain $B_L$ of $L$ copies of $H$.
Between consecutive copies, subdivide the right reserved edges of the earlier copy and the left reserved edges of the later copy.
Join the resulting $K$ ports by an arbitrary matching.
Each intermediate graph has cyclic connectivity at least $K$, by the first part of Lemma~\ref{lem:random-distant-port-gluing}.
The required port distances hold at every step.
A path between two new ports either stays within their end block or first reaches an earlier attachment port.
In either case the reserved-edge separation gives length at least $8K$.
Subdivisions cannot reduce these distances.

All edges of $B_L$ lie within a block or between consecutive blocks.
It has
\[
 |V(B_L)|=L n_0+2K(L-1)
\]
vertices.
Its girth is at least $K$.
Indeed, an old cycle remains at least as long, and a cycle crossing a newly added joining cut uses two distinct ports in each participating side.
The port-distance bound makes such a cycle longer than $K$.

Put
\[
 t=\left\lfloor\frac{L-1}{k-1}\right\rfloor.
\]
In blocks $1,1+t,\ldots,1+(k-1)t$, subdivide the extra reserved edge.
These $k$ new vertices form the ports of a normalized $k$-pole $X_L$.
Any path between different ports must traverse at least $t$ consecutive block boundaries.
Thus their pairwise distances are at least $t$.
The internal cyclic connectivity remains at least $K$, since subdivision preserves cyclic cuts and cycles.

For $L$ large enough that $t\ge2k$, join two disjoint copies of $X_L$ through their corresponding ports.
The second part of Lemma~\ref{lem:random-distant-port-gluing} makes this joining cut the unique minimum cyclic cut, of order $k$.
Any cycle meeting it uses at least two joining edges.
Between consecutive joining edges, it contains a path between distinct ports of one pole.
Hence its length is at least $2t+2$.
All other cycles lie in a copy of $X_L$ and have length at least $K$.

The rank of the final graph is
\[
 r=L(n_0+2K)-k-1.
\]
Since $k,n_0,K$ are fixed and $2t+2$ grows linearly with $L$, the claimed constant $c_k>0$ exists.
\end{proof}

In particular, no universal rule can select an edge of a minimum cyclic cut that lies on a polylogarithmic cycle.
The same construction rules out a uniform upper bound $r/(\log r)^2$ for that cycle length.
These statements already fail at each fixed $k\ge5$.
The graphs above are covered by the fixed-cut sharp bound, so this construction does not contradict that theorem.
It also leaves open a short-cycle selection rule under an additional near-extremal cycle-count hypothesis.

\begin{corollary}[The obstruction with growing connectivity and girth]
\label{cor:random-long-cut-diagonal}
There is a sequence of connected simple cubic graphs $G_k$ such that
\[
 \lambda_c(G_k)=k\longrightarrow\infty,
 \qquad
 \operatorname{girth}(G_k)\ge k+1\longrightarrow\infty,
\]
and every edge of their unique minimum cyclic cut lies on no cycle of length at most
\[
 \frac{r_k}{(\log r_k)^2},
 \qquad r_k=\beta(G_k).
\]
\end{corollary}

\begin{proof}
For the fixed-$k$ family in Theorem~\ref{thm:random-long-minimum-cut}, choose a rank $r_k$ so large that
\[
 c_k(\log r_k)^2>1.
\]
Every cycle through a minimum-cut edge then has length at least $c_k r_k>r_k/(\log r_k)^2$.
Choose the ranks increasing with $k$.
\end{proof}
